\documentclass{article}
\usepackage{mathnotes}

\title{Flows on the Line}

\begin{document}

We use $\dot{x}$ to represent the time derivative of $x,$ that is, $\dot{x} = dx/dt.$

\begin{definition}[First-order system]\synonyms{"one-dimensional system"}
A \textbf{first-order} system is a system of the form

\[ \dot{x} = f(x), \]

where $x = x(t)$ is a \@{real-valued} \@{function} of time $t,$ and $f(x)$ is a \@{smooth} real-valued function of $x.$
\end{definition}

We can interpret a \@{differential-equation} of this form as a as a \@{vector-field}. $f(x)$ tells us in which direction and at what magnitude we move at any point on the real line (where $f(x)$ is defined.) We can do this by plotting $\dot{x}$ vs $x.$

\includedemo{flows-on-the-line}

You can see that when $x$ is to the right of $0,$ the flow is to the right, and when $x$ is to the left of 0 the flow is to the left.

At some points, there is no flow.

\begin{definition}[fixed point]\synonyms{"equilibrium"}
A point at which the change in position of a system is $0$ is called a \textbf{fixed point.} We denote a fixed point at $x_p$ as $x^* = x_p.$
\end{definition}

\begin{definition}[Stable]
A \@{fixed-point} where nearby points flow towards it is said to be \textbf{stable.}
\end{definition}

\begin{definition}[Unstable]
A \@{fixed-point} where nearby points flow away from it is said to be \textbf{unstable.}
\end{definition}

\begin{remark}
The convention, when drawing a dynamical system as a vector field is to draw filled in circles for \@{stable} \@{fixed-points} and empty circles for \@{unstable} \@{fixed-points}.
\end{remark}

\begin{definition}[Phase point]
The starting point, $x_0,$ where we place a particle is called a \textbf{phase point.}
\end{definition}

\begin{definition}[Trajectory]
The function describing the path taken by a particle starting at a \@{phase-point} is called a \textbf{trajectory} and represents a solution to a \@{differential-equation} with \@{initial-conditions} $x = x_0.$
\end{definition}

\begin{definition}[Phase Portrait]
A drawing that shows the different \@{trajectories} taken from different \@{phase-points} in a system is called a \textbf{phase portrait.}
\end{definition}

We can also visualize trajectories by plotting $x(t)$ vs $t.$

\includedemo{time-evolution}

Here's both views together.

\includedemo{phase-and-time}

\begin{definition}[Globally stable]
A \@{fixed-point} $x^*$ that is approached from any starting position on the real line (other than that at $x^*$ itself) is said to be \textbf{globally stable.}
\end{definition}

\subsection{Application: Population Growth}

A very simple model of population growth is just exponential growth. You can model this as

\[ \dot{N} = r{N}, \]

where $r > 0.$ You can see by modeling this on the demo's above that population just goes to infinity with this. This is not realistic. A better model assumes there is a certain carrying capacity $K$, where if the population $N$ exceeds $K$, growth actually becomes negative. This is modeled using the logistic equation

\[ \dot{N} = rN(1 - \frac{N}{K}). \]

\subsection{Linear Stability Analysis}

While it's nice to have a visual intuition for whether a \@{fixed-point} is stable or not, sometimes it's also nice to know analytically.

\begin{theorem}\label{linear-stability-analysis}
Let $x^*$ be a fixed point of $\dot{x} = f(x).$ Then, if $f'(x) \neq 0,$ if $f'(x^*)$ is negative, then $x^*$ is a \@{stable} \@{fixed-point}
. If $f'(x^*)$ is positive, then $x^*$ is an \@{unstable} \@{fixed-point}.
\end{theorem}

\begin{remark}
This comes from letting $u(t) = x(t) - x^*$ be a small perturbation away from $x^*,$ differentiating it, writing its \@{taylor-series}, then noticing that $f(x^*) = 0$ and terms greater than the linear term matter less than the linear term and writing

\[ \dot{u} = f'(x^*)u. \]

This is called  the \textbf{linearization} about $x^*.$

It also only works if $f'(x) \neq 0.$ If that's not the case, the best bet is to fall back to graphical analysis or to solve explicitly if possible.
\end{remark}

\begin{definition}[Half-stable]
A \@{fixed-point} $x^*$ where $f'(x) = 0$ and the stability depends on which side of $x^*$ $x_0$ lies on is called \textbf{half-stable.} It is denoted with a half-filled circle.
\end{definition}

\subsection{Existence and Uniqueness}

\begin{theorem}\label{existence-and-uniqueness-of-ivp-solutions}
Consider the \@{initial-value-problem}

\[ \dot{x} = f(x), \quad x(0) = x_0. \]

Suppose that $f(x)$ and $f'(x)$ are \@{continuous} on an \@{open} \@{interval} $R$ of the $x$-axis, and suppose that $x_0$ is a \@{point} in $R.$ Then, the initial value problem has a solution $x(t)$ on some time interval $(- \tau, \tau)$ about $t = 0,$ and the solution is unique.
\end{theorem}

\end{document}
